\documentclass[11pt,a4paper]{article}

\usepackage[utf8]{inputenc}
\usepackage[T1]{fontenc}
\usepackage[french]{babel}

\usepackage{styles/couleurs}
\usepackage{styles/mise_en_page}
\usepackage{styles/boites}
\usepackage{styles/maths}
\usepackage{styles/figures}

\usepackage{array}
\usepackage{tabularx}
\usepackage{pgffor}
\usepackage{pdflscape}

\geometry{
  top=1.1cm,
  bottom=1.1cm,
  left=1.4cm,
  right=1.4cm
}

\setlength{\parindent}{0pt}
\setlength{\parskip}{0.15em}

\definecolor{DefiC}{HTML}{8B4A97}
\definecolor{DefiBg}{HTML}{F7F0F8}

% =========================================================
% MACROS
% =========================================================

\newcommand{\lignes}[1]{%
\foreach \i in {1,...,#1}{%
  \noindent\dotfill\par\vspace{0.45em}
}
}

\newcommand{\stopprof}{%
\begin{center}
\begin{tcolorbox}[
  width=5.8cm,
  colback=ActionC!6,
  colframe=ActionC!70,
  boxrule=0.8pt,
  arc=2mm,
  halign=center,
  top=1.5mm,
  bottom=1.5mm,
  before skip=1.5mm,
  after skip=1.5mm
]
{\sffamily\bfseries\color{ActionC}
STOP PROF
\qquad
$\square$}
\end{tcolorbox}
\end{center}
}

% =========================================================
% CARTES DE BASE
% =========================================================

\newtcolorbox{cartebase}[1]{
  enhanced,
  colback=white,
  colframe=NeutreC!65,
  boxrule=0.7pt,
  arc=1mm,
  height=4.15cm,
  valign=center,
  left=3mm,
  right=3mm,
  top=2mm,
  bottom=2mm,
  title={\sffamily\bfseries Carte #1},
  coltitle=TitreC,
  colbacktitle=NeutreBg
}

\newcommand{\carte}[2]{%
\begin{minipage}[t]{0.242\linewidth}
\begin{cartebase}{#1}
\centering
#2
\end{cartebase}
\end{minipage}
}

% =========================================================
% CARTES DÉFI
% =========================================================

\newtcolorbox{cartedefi}[1]{
  enhanced,
  colback=white,
  colframe=DefiC!70,
  boxrule=0.8pt,
  arc=1mm,
  height=7.2cm,
  valign=center,
  left=4mm,
  right=4mm,
  top=3mm,
  bottom=3mm,
  title={\sffamily\bfseries Défi #1},
  coltitle=DefiC,
  colbacktitle=DefiBg
}

\newcommand{\defi}[2]{%
\begin{minipage}[t]{0.242\linewidth}
\begin{cartedefi}{#1}
\centering
#2
\end{cartedefi}
\end{minipage}
}

% =========================================================
% DOCUMENT
% =========================================================

\begin{document}

% =========================================================
% PAGE 1 -- REGROUPER + CALCULER
% =========================================================

\begin{center}
{\sffamily\LARGE\bfseries\color{TitreC}
AP -- Calcul réfléchi -- Sans calculatrice}
\end{center}

\vspace{0.3em}

% ---------------------------------------------------------
% 1. REGROUPER
% ---------------------------------------------------------

\begin{tcolorbox}[
  colback=white,
  colframe=SavoirC!60,
  boxrule=0.7pt,
  arc=2mm,
  left=4mm,
  right=4mm,
  top=2.5mm,
  bottom=2.5mm,
  before skip=1mm,
  after skip=1mm
]

{\sffamily\Large\bfseries\color{SavoirC}
1. Regrouper}

\smallskip

Regroupez les cartes qui, selon vous, peuvent se résoudre avec
\textbf{la même idée de calcul}.

\smallskip

En utilisant vos mots, donnez un nom à chaque groupe.

\medskip

\textbf{Noms de nos groupes :}

\vspace{0.25em}

\begin{tcolorbox}[
  colback=NeutreBg!20,
  colframe=NeutreC!30,
  boxrule=0.4pt,
  arc=1.5mm,
  height=3.6cm,
  left=3mm,
  right=3mm,
  top=2mm,
  bottom=2mm
]
\end{tcolorbox}

\end{tcolorbox}

\stopprof

% ---------------------------------------------------------
% 2. CALCULER
% ---------------------------------------------------------

\begin{tcolorbox}[
  colback=white,
  colframe=SavoirC!60,
  boxrule=0.7pt,
  arc=2mm,
  left=4mm,
  right=4mm,
  top=2.5mm,
  bottom=2.5mm,
  before skip=1mm,
  after skip=1mm
]

{\sffamily\Large\bfseries\color{SavoirC}
2. Calculer}

\smallskip

Choisissez \textbf{4 cartes de groupes différents}.

Résolvez-les en expliquant votre raisonnement.

\medskip

\begin{center}

\renewcommand{\arraystretch}{1}

\begin{tabularx}{0.98\linewidth}{
|>{\centering\arraybackslash}p{1.7cm}|
>{\raggedright\arraybackslash}X|
}

\hline
\textbf{Carte}
&
\textbf{Notre raisonnement}
\\
\hline

\rule{0pt}{2.25cm}
&
\\
\hline

\rule{0pt}{2.25cm}
&
\\
\hline

\rule{0pt}{2.25cm}
&
\\
\hline

\rule{0pt}{2.25cm}
&
\\
\hline

\end{tabularx}

\end{center}

\end{tcolorbox}

\stopprof

% =========================================================
% PAGE 2 -- DEUX CHEMINS + CONTRÔLER
% =========================================================

\newpage

% ---------------------------------------------------------
% 3. DEUX CHEMINS
% ---------------------------------------------------------

\begin{tcolorbox}[
  colback=white,
  colframe=SavoirC!60,
  boxrule=0.7pt,
  arc=2mm,
  left=4mm,
  right=4mm,
  top=3mm,
  bottom=3mm
]

{\sffamily\Large\bfseries\color{SavoirC}
3. Deux chemins}

\medskip

Choisissez \textbf{une carte}.

Trouvez \textbf{deux façons différentes} de la résoudre.

\medskip

\textbf{Carte :}
\rule{2cm}{0.4pt}

\vspace{0.7em}

\textbf{Chemin 1}

\vspace{0.2em}

\lignes{6}

\vspace{0.6em}

\textbf{Chemin 2}

\vspace{0.2em}

\lignes{6}

\vspace{0.6em}

\textbf{Lequel préférez-vous ? Pourquoi ?}

\vspace{0.2em}

\lignes{3}

\end{tcolorbox}

\vspace{0.5em}

% ---------------------------------------------------------
% 4. CONTRÔLER
% ---------------------------------------------------------

\begin{tcolorbox}[
  colback=white,
  colframe=SavoirC!60,
  boxrule=0.7pt,
  arc=2mm,
  left=4mm,
  right=4mm,
  top=3mm,
  bottom=3mm
]

{\sffamily\Large\bfseries\color{SavoirC}
4. Contrôler}

\medskip

Choisissez un résultat.

Contrôlez-le \textbf{sans refaire le même calcul}.

\vspace{0.4em}

\lignes{7}

\end{tcolorbox}

\stopprof

% =========================================================
% PAGE 3 -- TRACE ÉCRITE
% =========================================================

\newpage

\begin{center}
{\sffamily\LARGE\bfseries\color{TitreC}
Nos outils pour calculer}
\end{center}

\vspace{0.8em}

\textit{À compléter ensemble.}

\vspace{1em}

\begin{center}

\renewcommand{\arraystretch}{1}

\begin{tabularx}{0.96\linewidth}{
|>{\raggedright\arraybackslash}p{10.5cm}|
>{\raggedright\arraybackslash}X|
}

\hline

\textbf{Ce qu'on peut faire}
&
\textbf{Un exemple}
\\
\hline

\rule{0pt}{4cm}
&
\\
\hline

\rule{0pt}{4cm}
&
\\
\hline

\rule{0pt}{4cm}
&
\\
\hline

\rule{0pt}{4cm}
&
\\
\hline

\rule{0pt}{4cm}
&
\\
\hline

\end{tabularx}

\end{center}

% =========================================================
% PAGE 4 -- CARTES DE BASE
% =========================================================

\clearpage

\newgeometry{margin=0.7cm}

\begin{landscape}

\pagestyle{empty}

\begin{center}
{\sffamily\Large\bfseries\color{TitreC}
Cartes}
\end{center}

\vspace{0.15cm}

% ---------------------------------------------------------
% LIGNE 1
% ---------------------------------------------------------

\noindent
\carte{A}{
{\Large
\[
48+39
\]
}
}
\hfill
\carte{B}{
{\Large
\[
25\times36
\]
}
}
\hfill
\carte{C}{
{\Large
\[
48\times5
\]
}
}
\hfill
\carte{D}{
{\Large
\[
125\times8
\]
}
}

\vspace{0.1cm}

% ---------------------------------------------------------
% LIGNE 2
% ---------------------------------------------------------

\noindent
\carte{E}{
{\Large
\[
302-198
\]
}
}
\hfill
\carte{F}{
{\Large
\[
99\times17
\]
}
}
\hfill
\carte{G}{
{\Large
\[
360\div5
\]
}
}
\hfill
\carte{H}{
{\Large
\[
8\times\boxed{\phantom{0}}=6{,}4
\]
}
}

\vspace{0.1cm}

% ---------------------------------------------------------
% LIGNE 3
% ---------------------------------------------------------

\noindent
\carte{I}{
{\Large
\[
\boxed{\phantom{00}}\div6=8
\]
}
}
\hfill
\carte{J}{
{\Large
\[
35\div\boxed{\phantom{0}}=7
\]
}
}
\hfill
\carte{K}{
{\Large
\[
\dfrac34\text{ de }28
\]
}
}
\hfill
\carte{L}{
{\Large
\[
25\,\%\text{ de }36
\]
}
}

\vspace{0.1cm}

% ---------------------------------------------------------
% LIGNE 4
% ---------------------------------------------------------

\noindent
\carte{M}{
{\Large
\[
0{,}5\times48
\]
}
}
\hfill
\carte{N}{
{\Large
\[
12\times7+12\times3
\]
}
}
\hfill
\carte{O}{
\large

\[
19{,}8\times5{,}1
\]

\vspace{0.2em}

Plus proche de

\[
10,\qquad100\qquad\text{ou}\qquad1000\ ?
\]
}
\hfill
\carte{P}{
\large

\[
398\div8=497{,}5
\]

\vspace{0.4em}

\textbf{Possible ?}

\vspace{0.2em}

Justifie.
}

\end{landscape}

% =========================================================
% PAGE 5 -- DÉFIS
% =========================================================

\begin{landscape}

\begin{center}

{\sffamily\Large\bfseries\color{DefiC}
Défi -- Deux idées à associer}

\vspace{0.25em}

{\large
Choisissez un défi.

Associez \textbf{deux idées de calcul}.

Nommez-les et résolvez.}

\end{center}

\vspace{0.55cm}

% ---------------------------------------------------------
% LIGNE 1
% ---------------------------------------------------------

\noindent
\defi{1}{
{\Large
\[
25\,\%\text{ de }\boxed{\phantom{00}}=18
\]
}
}
\hfill
\defi{2}{
{\Large
\[
\boxed{\phantom{00}}\div5
=
\dfrac34\text{ de }20
\]
}
}
\hfill
\defi{3}{
{\Large
\[
25\times48+25\times52
\]
}
}
\hfill
\defi{4}{
{\Large
\[
125\times24
\]
}
}

\vspace{0.45cm}

% ---------------------------------------------------------
% LIGNE 2
% ---------------------------------------------------------

\noindent
\defi{5}{
{\Large
\[
48\times2{,}5
\]
}
}
\hfill
\defi{6}{
{\Large
\[
0{,}75\times48
\]
}
}
\hfill
\defi{7}{
\large

\[
49{,}8\times0{,}21
\]

\vspace{0.3em}

Plus proche de

\[
1,\qquad10\qquad\text{ou}\qquad100\ ?
\]
}
\hfill
\defi{8}{
\large

\[
25\,\%\text{ de }84=42
\]

\vspace{0.5em}

\textbf{Possible ?}

\vspace{0.2em}

Justifie.
}

\end{landscape}

\restoregeometry

% =========================================================
% NOTES PROF -- NON IMPRIMÉES
% =========================================================
%
% ---------------------------------------------------------
% PASSAGE PROF 1
% ---------------------------------------------------------
%
% Vérifier qu'il existe une logique mathématique
% dans les regroupements.
%
% Il n'y a pas de classement unique attendu.
%
% Relances possibles :
% - Pourquoi ces deux cartes vont-elles ensemble ?
% - Cette carte pourrait-elle aller ailleurs ?
% - Quelle carte vous a fait hésiter ?
%
% ---------------------------------------------------------
% PASSAGE PROF 2
% ---------------------------------------------------------
%
% Choisir soi-même l'une des quatre cartes.
%
% Faire expliquer le raisonnement.
%
% Si l'idée est bonne mais le résultat faux :
% "Ton idée me va. Vérifie ton résultat."
%
% ---------------------------------------------------------
% PASSAGE PROF 3
% ---------------------------------------------------------
%
% Vérifier :
% - deux chemins réellement différents ;
% - un contrôle qui ne refait pas exactement
%   le calcul de départ.
%
% ---------------------------------------------------------
% IDÉES QUE LA MISE EN COMMUN PEUT FAIRE ÉMERGER
% ---------------------------------------------------------
%
% Ne pas donner cette liste aux élèves avant la trace.
%
% - Décomposer / recomposer
% - S'appuyer sur un calcul simple que l'on connaît
% - Utiliser une opération inverse
% - Changer d'écriture
% - Estimer / contrôler
%
% ---------------------------------------------------------
% CARTES DE BASE : PISTES
% ---------------------------------------------------------
%
% A
% 48 + 39
% = 48 + 40 - 1
% = 87
%
% ---------------------------------------------------------
%
% B
% 25 x 36
%
% 36 = 9 x 4
%
% 25 x 36
% = 25 x 4 x 9
% = 100 x 9
% = 900
%
% ---------------------------------------------------------
%
% C
% 48 x 5
% = 48 x 10 / 2
% = 480 / 2
% = 240
%
% ---------------------------------------------------------
%
% D
% 125 x 8
% = 1000
%
% ---------------------------------------------------------
%
% E
% 302 - 198
% = 302 - 200 + 2
% = 104
%
% Autre possibilité :
% chercher l'écart entre 198 et 302.
%
% ---------------------------------------------------------
%
% F
% 99 x 17
% = 100 x 17 - 17
% = 1683
%
% ---------------------------------------------------------
%
% G
% 360 / 5
%
% 360 / 5
% = 720 / 10
% = 72
%
% ---------------------------------------------------------
%
% H
% 8 x ? = 6,4
% ? = 0,8
%
% ---------------------------------------------------------
%
% I
% ? / 6 = 8
% ? = 48
%
% ---------------------------------------------------------
%
% J
% 35 / ? = 7
% ? = 5
%
% ---------------------------------------------------------
%
% K
% 3/4 de 28
% = 21
%
% ---------------------------------------------------------
%
% L
% 25 % de 36
% = 1/4 de 36
% = 9
%
% ---------------------------------------------------------
%
% M
% 0,5 x 48
% = moitié de 48
% = 24
%
% ---------------------------------------------------------
%
% N
% 12 x 7 + 12 x 3
% = 12 x (7 + 3)
% = 120
%
% ---------------------------------------------------------
%
% O
% 19,8 x 5,1
% environ 20 x 5
% donc proche de 100
%
% ---------------------------------------------------------
%
% P
% 398 / 8
% environ 400 / 8 = 50
%
% 497,5 est impossible.
%
% =========================================================
% DÉFIS
% =========================================================
%
% DÉFI 1
%
% 25 % de ? = 18
%
% 25 % = 1/4
%
% Si le quart vaut 18,
% le tout vaut 72.
%
% ---------------------------------------------------------
%
% DÉFI 2
%
% ? / 5 = 3/4 de 20
%
% 3/4 de 20 = 15
%
% ? / 5 = 15
%
% ? = 75
%
% ---------------------------------------------------------
%
% DÉFI 3
%
% 25 x 48 + 25 x 52
%
% = 25 x (48 + 52)
% = 25 x 100
% = 2500
%
% ---------------------------------------------------------
%
% DÉFI 4
%
% 125 x 24
%
% 24 = 8 x 3
%
% 125 x 24
% = 125 x 8 x 3
% = 1000 x 3
% = 3000
%
% ---------------------------------------------------------
%
% DÉFI 5
%
% 48 x 2,5
%
% 2,5 = 5 / 2
%
% 48 x 2,5
% = 48 x 5 / 2
% = 120
%
% Autre possibilité :
%
% 2,5 = 2 + 0,5
%
% 48 x 2 + moitié de 48
% = 96 + 24
% = 120
%
% ---------------------------------------------------------
%
% DÉFI 6
%
% 0,75 x 48
%
% 0,75 = 3/4
%
% 3/4 de 48 = 36
%
% ---------------------------------------------------------
%
% DÉFI 7
%
% 49,8 x 0,21
%
% environ 50 x 0,2
% donc environ 10
%
% ---------------------------------------------------------
%
% DÉFI 8
%
% 25 % de 84 = 42
%
% 25 % = 1/4
%
% 84 / 4 = 21
%
% donc 42 est faux.
%
% ---------------------------------------------------------
%
% Les défis sont une réserve pour les groupes rapides.
% Il n'est pas demandé de finir le paquet.

\end{document}